\documentclass[10pt]{article}
\usepackage[T1]{fontenc}
\usepackage[utf8]{inputenc}
\usepackage[french]{babel}
\usepackage{amsmath,amssymb,mathrsfs,mathpazo}
\usepackage{geometry}
\geometry{a4paper,margin=2cm}
\usepackage{graphicx,array,tabularx,tabularray,colortbl,multirow,multicol,booktabs,enumitem,pifont,ulem,fancybox,diagbox,cancel}
\usepackage{pstricks,pstricks-add,pst-plot,pst-tree,pst-node,pst-text,pst-func,pst-math,pst-bspline,pst-3dplot,pst-eucl,pst-solides3d}
\let\pscalculate\undefined
\usepackage{pst-bezier}
\usepackage{tikz,pgfplots,tkz-tab}
\usetikzlibrary{arrows,arrows.meta,patterns,calc,positioning,trees,decorations.markings,shapes.geometric,shapes.misc,shapes.arrows,intersections,decorations.pathmorphing,decorations.pathreplacing,backgrounds,shadows,babel}
\pgfplotsset{compat=1.18}
\usepgfplotslibrary{fillbetween}
\usepackage[np]{numprint}
\usepackage{listings,makecell,pas-tableur,twoopt}
\let\barre\undefined
\usepackage[pstricks]{tablvar}
\definecolor{darkgreen}{RGB}{0,100,0}
\providecommand{\red}{\color{red}}
\providecommand{\blue}{\color{blue}}
\providecommand{\green}{\color{green}}
\providecommand{\black}{\color{black}}
\providecommand{\white}{\color{white}}
\providecommand{\gray}{\color{gray}}
\providecommand{\cyan}{\color{cyan}}
\providecommand{\magenta}{\color{magenta}}
\providecommand{\yellow}{\color{yellow}}
\providecommand{\violet}{\color{violet}}
\providecommand{\orange}{\color{orange}}
\providecommand{\euro}{\texteuro}
\providecommand{\EUR}{\texteuro}
\providecommand{\ieme}{\textsuperscript{e}}
\providecommand{\degres}{\ensuremath{^\circ}}
\providecommand{\degree}{\ensuremath{^\circ}}
\providecommand{\vv}[1]{\overrightarrow{#1}}
\providecommand{\Coord}[1]{#1}
\providecommand{\decosix}{}
\providecommand{\starredbullet}{\ensuremath{\bullet}}
\pagestyle{empty}
\setlength{\parindent}{0pt}
\usepackage{adjustbox}\DeclareUnicodeCharacter{2236}{:}\DeclareUnicodeCharacter{2190}{\ensuremath{\leftarrow}}
\providecommand{\R}{\mathbb{R}}
\providecommand{\N}{\mathbb{N}}
\providecommand{\Z}{\mathbb{Z}}
\providecommand{\Q}{\mathbb{Q}}
\providecommand{\D}{\mathbb{D}}
\providecommand{\C}{\mathbb{C}}
\providecommand{\e}{{\mathrm{e}}}
\providecommand{\vect}[1]{\overrightarrow{#1}}
\providecommand{\vectt}[1]{\overrightarrow{\mathrm{#1}}}
\providecommand{\vvt}[1]{\overrightarrow{\mathrm{#1}}}
\providecommand{\np}[2][]{{#2~#1}}
\providecommand{\vv}[1]{{\overrightarrow{#1}}}
\renewcommand{\d}{{\mathrm{d}}}
\providecommand{\degres}{{^\circ}}
\providecommand{\degre}{{^\circ}}
\providecommand{\blue}{\color{blue}}
\providecommand{\red}{\color{red}}
\providecommand{\enskip}{\quad}
\providecommand{\llbracket}{\lbrack\!\lbrack}
\providecommand{\rrbracket}{\rbrack\!\rbrack}
\providecommand{\euro}{{\text{€}}}
\providecommand{\no}{{\text{n°}}}
\providecommand{\interval}[2]{{[#1;#2]}}
\providecommand{\intervFO}[2]{{[#1;#2[}}
\providecommand{\dotfill}{\ldots}

\providecommand{\ds}{\displaystyle}
\providecommand{\textding}[1]{#1}
\providecommand{\Oij}{\(\mathrm{O};\vec{\imath},\vec{\jmath}\)}
\providecommand{\Oijk}{\(\mathrm{O};\vec{\imath},\vec{\jmath},\vec{k}\)}
\providecommand{\Ouv}{\(\mathrm{O};\vec{u},\vec{v}\)}
\providecommand{\intervFF}[2]{\ensuremath{[#1;#2]}}
\renewcommand{\euro}{\ifmmode\text{\texteuro}\else\texteuro\fi}
\DeclareUnicodeCharacter{00B0}{\ifmmode^{\circ}\else\textdegree\fi}
\DeclareUnicodeCharacter{20AC}{\ifmmode\text{\texteuro}\else\texteuro\fi}
% Énoncés : Éducation nationale. Transcription et corrigés : APMEP, Denis Vergès et contributeurs. Découpage et classement : CoopMaths, Sébastien Lozano, Éric Elter et contributeurs. Adaptation de la présentation : Maths Explorer.
% CC BY-SA : https://coopmaths.fr/www/about/
% https://coopmaths.fr/alea/static/e3c/2020/tex/e3c_2020_00_sujet18_3.tex
\begin{document}
\noindent\textit{APMEP / CoopMaths. CC BY-SA. Présentation adaptée par Maths Explorer.}\par\medskip
\medskip

\emph{Dans cet exercice toutes les probabilités seront données sous forme décimale, arrondie au
millième.}
\medskip


Une entreprise récupère des smartphones endommagés, les répare et les reconditionne afin de
les revendre à prix réduit.

\begin{itemize}
\item  45\,\% des smartphones qu’elle récupère ont un écran cassé ; 
\item  parmi les smartphones ayant un écran cassé, 30\,\% ont également une batterie
défectueuse ;
 \item par contre, seulement 20\,\% des smartphones ayant un écran non cassé ont une batterie
défectueuse.
\end{itemize}
\begin{enumerate}
\item Un technicien chargé de réparer et reconditionner les smartphones de l’entreprise prend un
smartphone au hasard dans le stock. On note :

\begin{itemize}
\item $E$ l’évènement : \og Le smartphone choisi a un écran cassé \fg.
\item $B$ l’évènement : \og Le smartphone choisi a une batterie défectueuse \fg.
\end{itemize}

\medskip

\begin{enumerate}
\item Représenter la situation décrite ci-dessus par un arbre pondéré.
\item Démontrer que la probabilité que le smartphone choisi ait une batterie défectueuse est
égale à $0,245$.
\item Sachant que le smartphone choisi a une batterie défectueuse, quelle est la probabilité
qu’il ait un écran cassé ?
\end{enumerate}
\item L’entreprise dépense 20 \euro{} pour réparer et reconditionner chaque smartphone qu’elle
 récupère. Si l’écran est cassé, elle dépense 30 \euro{} supplémentaires, et si la batterie est
défectueuse, elle dépense 40 \euro{} supplémentaires.

On note $X$ la variable aléatoire égale au coût total de réparation et reconditionnement d’un
smartphone choisi au hasard dans le stock.
\begin{enumerate}
\item Recopier et compléter sur la copie (aucune justification n’est attendue) le tableau
suivant pour donner la loi de probabilité de la variable aléatoire $X$.

\begin{center}
\begin{tabularx}{\linewidth}{|*{5}{>{\centering \arraybackslash} X|}} \hline
 $x_i$		&20		&50			&$\cdots$	&$\cdots$\\ \hline
 $p(X=x_i)$	&0,44	&$\cdots$  	&$\cdots$ 	&$\cdots$ \\ \hline
\end{tabularx}
\end{center}
\item L’entreprise doit réparer et reconditionner $500$ smartphones. Combien doit-elle
s’attendre à dépenser ?
\end{enumerate}
\end{enumerate}

\vspace{0,5cm}

\end{document}
